\chapter{The Horizon in Flat Coordinates and the Thermodynamics of Evaporation}

\lectureoverview{The Schwarzschild metric appears to become singular at the horizon, although a freely falling observer should notice nothing exceptional there. We resolve that tension by converting the near-horizon metric to flat Minkowski space in hyperbolic polar coordinates. The same geometry then separates Alice's infalling description from Bob's accelerated exterior description. The second half develops entropy and temperature, derives Hawking's inverse-mass temperature and the \(M^3\) evaporation law, and closes by asking what information conservation requires of the microscopic theory of gravity.}

\section{The Puzzle at \texorpdfstring{\(r=2GM\)}{r=2GM}}

Begin with a very large black hole. Its horizon can have such a small curvature that an observer falling through it sees, over any modest laboratory-sized region, nearly empty flat spacetime. No wall is struck, and no local alarm announces the crossing. Yet the Schwarzschild metric seems to say that something singular occurs at precisely that radius. The purpose of the first half of the lecture is to make these two statements compatible.

We use the mostly-plus convention. In flat spacetime,
\begin{equation}
ds^2=-dt^2+dx^2+dy^2+dz^2,
\end{equation}
so a timelike worldline has \(ds^2<0\) and proper time satisfies \(d\tau^2=-ds^2\). This is the overall sign reversal of the convention used in the preceding lecture; no observable depends on that choice.

With \(c=1\), the Schwarzschild line element is
\begin{equation}
\begin{aligned}
ds^2={}&-\left(1-\frac{2GM}{r}\right)dt^2
+\frac{dr^2}{1-\frac{2GM}{r}}
+r^2d\Omega_2^2,\\
d\Omega_2^2={}&d\theta^2+\sin^2\theta\,d\phi^2 .
\end{aligned}
\label{eq:l03-schwarzschild}
\end{equation}
The Schwarzschild radius is \(R_s=2GM\), or \(R_s=2GM/c^2\) when \(c\) is restored. Far away, \(r\gg R_s\), both radial coefficients approach their flat-space values. At \(r=R_s\), however, the coefficient of \(dt^2\) vanishes and the coefficient of \(dr^2\) diverges. That looks dramatic. The question is whether the geometry is dramatic or only the coordinates.

\section{Two Coordinate Lessons}

\subsection{The harmless origin of polar coordinates}

On an ordinary Euclidean plane,
\begin{equation}
ds^2=dx^2+dy^2.
\end{equation}
Introduce polar coordinates
\begin{equation}
x=\rho\cos\theta,
\qquad
y=\rho\sin\theta .
\end{equation}
A small displacement has radial length \(d\rho\) and transverse length \(\rho\,d\theta\), hence
\begin{equation}
ds^2=d\rho^2+\rho^2d\theta^2.
\label{eq:l03-polar}
\end{equation}
At \(\rho=0\), the angular coefficient vanishes. This does not make the center of the plane singular. It says only that every value of \(\theta\) names the same point there. A peculiar metric coefficient may therefore diagnose a peculiar coordinate system rather than peculiar physics.

\subsection{Hyperbolic polar coordinates}

Now replace the Euclidean plane by \(1+1\)-dimensional Minkowski spacetime,
\begin{equation}
ds^2=-dT^2+dX^2.
\end{equation}
In the right-hand wedge \(X>|T|\), define
\begin{equation}
X=\rho\cosh\omega,
\qquad
T=\rho\sinh\omega,
\qquad \rho>0 .
\label{eq:l03-rindler-transform}
\end{equation}
The ordinary identity \(\cos^2\theta+\sin^2\theta=1\) is replaced by
\begin{equation}
\cosh^2\omega-\sinh^2\omega=1.
\end{equation}
Thus a line of constant \(\rho\) is the hyperbola
\begin{equation}
X^2-T^2=\rho^2.
\end{equation}
Unlike the periodic angle \(\theta\), the rapidity-like coordinate \(\omega\) runs from \(-\infty\) to \(+\infty\). At large positive \(\omega\), \(\cosh\omega\) and \(\sinh\omega\) become nearly equal and the hyperbola approaches the future null ray \(X=T\).

Differentiating Equation~\eqref{eq:l03-rindler-transform} gives
\begin{align}
dX&=\cosh\omega\,d\rho+\rho\sinh\omega\,d\omega,\\
dT&=\sinh\omega\,d\rho+\rho\cosh\omega\,d\omega .
\end{align}
The cross terms cancel, leaving
\begin{equation}
ds^2=-\rho^2d\omega^2+d\rho^2.
\label{eq:l03-rindler-metric}
\end{equation}

\begin{figure}[htbp]
\centering
\includegraphics[width=0.72\textwidth]{lecture_03_figure_02.png}
\caption{The completed hyperbolic-coordinate board: \(X=\rho\cosh\omega\), \(T=\rho\sinh\omega\), the unit-hyperbola construction, and the flat metric in the form \(ds^2=-\rho^2d\omega^2+d\rho^2\).}
\lectureframe{00:16:55.000}
\end{figure}

Equation~\eqref{eq:l03-rindler-metric} still describes perfectly flat spacetime. Its time coefficient vanishes at \(\rho=0\), just as the angular coefficient of Equation~\eqref{eq:l03-polar} vanishes at the origin. The coordinate chart degenerates there; spacetime does not.

\section{Zooming In on the Schwarzschild Horizon}

Return to the radial-time part of Equation~\eqref{eq:l03-schwarzschild}:
\begin{align}
ds^2
&=-\frac{r-2GM}{r}\,dt^2
+\frac{r}{r-2GM}\,dr^2 .
\label{eq:l03-rewritten-schwarzschild}
\end{align}
We now inspect only a narrow neighborhood of \(r=2GM\). In factors that vary slowly, \(r\) may be replaced by \(2GM\). The factor \(r-2GM\) must be retained, because it vanishes and changes sign at the horizon. Therefore
\begin{equation}
ds^2\approx
-\frac{r-2GM}{2GM}\,dt^2
+\frac{2GM}{r-2GM}\,dr^2 .
\label{eq:l03-near-horizon-r}
\end{equation}

Choose a new radial coordinate \(\rho\) so that the second term is simply \(d\rho^2\):
\begin{equation}
d\rho^2=\frac{2GM}{r-2GM}\,dr^2.
\end{equation}
Outside the horizon, taking the positive square root and integrating from \(r=2GM\) gives
\begin{align}
d\rho&=\sqrt{\frac{2GM}{r-2GM}}\,dr,\\
\rho&=2\sqrt{2GM}\sqrt{r-2GM},\\
\rho^2&=8GM(r-2GM),
\qquad
r-2GM=\frac{\rho^2}{8GM}.
\label{eq:l03-rho-r}
\end{align}
Substitution into Equation~\eqref{eq:l03-near-horizon-r} yields
\begin{equation}
ds^2\approx
-\frac{\rho^2}{16G^2M^2}\,dt^2+d\rho^2.
\end{equation}
Finally rescale Schwarzschild time,
\begin{equation}
\omega=\frac{t}{4GM},
\qquad
d\omega^2=\frac{dt^2}{16G^2M^2},
\label{eq:l03-omega}
\end{equation}
and obtain
\begin{equation}
\boxed{ds^2\approx-\rho^2d\omega^2+d\rho^2.}
\label{eq:l03-near-horizon-rindler}
\end{equation}

\begin{figure}[htbp]
\centering
\includegraphics[width=0.70\textwidth]{lecture_03_figure_03.png}
\caption{The board after the time rescaling: the near-horizon Schwarzschild metric has become the same hyperbolic-polar metric as flat Minkowski spacetime.}
\lectureframe{00:31:35.000}
\end{figure}

The approximation improves without limit as the inspected neighborhood shrinks toward the horizon. It proves the local statement that was wanted: for a sufficiently large hole, the horizon is ordinary spacetime written in an extraordinary chart. Curvature is not literally zero for a finite-mass Schwarzschild hole, but it is finite at the horizon and can be made arbitrarily small there by increasing \(M\).\editorialnote{The local curvature scale is of order \(R_s^{-2}\). The Rindler metric is the leading term in a local near-horizon expansion, not a claim that the entire Schwarzschild geometry is globally flat.}

\begin{classroomqa}
\audiencequestion{For \(r<2GM\), Equation~\eqref{eq:l03-rho-r} makes this exterior \(\rho\) imaginary. Does that signal a physical problem?\lecturetimestamp{00:32:30}}
\lecturerresponse{It reproduces the same coordinate failure in a different form. The present chart was constructed for the exterior wedge; extending through the horizon requires another chart. The failure of this particular \(\rho\) is not a physical event at the horizon.}
\end{classroomqa}

\begin{classroomqa}
\audiencequestion{In the near-horizon rewrite, how does the approximated line follow from the preceding exact line?\lecturetimestamp{00:39:21}}
\lecturerresponse{Only the slowly varying occurrences of \(r\) have been set equal to \(2GM\). The small difference \(r-2GM\) is left intact. That is the approximation, and it becomes increasingly accurate as \(r\) approaches \(2GM\).}
\end{classroomqa}

\section{What the Coordinates Mean}

The transformation identifies
\begin{equation}
\rho=0\quad\Longleftrightarrow\quad r=2GM.
\end{equation}
The exterior \(r>2GM\) is the right Rindler wedge \(X>|T|\). Constant-\(\rho\) hyperbolae represent observers who remain at fixed Schwarzschild radius. An outgoing light ray emitted anywhere in that wedge crosses every larger-\(\rho\) worldline, so an exterior observer can communicate outward. A point beyond the future null boundary cannot send a light signal back into the right wedge. That null boundary is the horizon.

\begin{classroomqa}
\audiencequestion{Does the hyperbolic-coordinate patch shown here always describe the region outside the horizon?\lecturetimestamp{00:35:22}}
\lecturerresponse{Yes. In this construction \(r>2GM\) gives \(\rho^2>0\) and places the event in the right exterior wedge. The null edge \(\rho=0\) is the horizon; the present coordinates do not continue as real exterior coordinates beyond it.}
\end{classroomqa}

\begin{classroomqa}
\audiencequestion{Does a constant-\(\rho\) hyperbola represent one time or one distance?\lecturetimestamp{00:37:21}}
\lecturerresponse{It represents one fixed proper distance from the horizon in the near-horizon approximation, and therefore one fixed Schwarzschild radius. Motion along the hyperbola is the passage of time for the observer stationed there.}
\end{classroomqa}

Along a fixed-\(\rho\) trajectory, Equation~\eqref{eq:l03-rindler-metric} gives
\begin{equation}
d\tau=\rho\,d\omega.
\label{eq:l03-proper-time-rindler}
\end{equation}
Thus \(\omega\) is not itself the observer's proper time. It is an angular time common to the family of hyperbolae; each observer multiplies it by that observer's \(\rho\), just as an ordinary angle is multiplied by a circle's radius to give arc length.

\begin{classroomqa}
\audiencequestion{Is \(\omega\) the proper time of the accelerated observer?\lecturetimestamp{00:51:45}}
\lecturerresponse{Almost. At fixed \(\rho\), the proper-time increment is \(\rho\,d\omega\). The factor of \(\rho\) is essential.}
\end{classroomqa}

\begin{figure}[htbp]
\centering
\includegraphics[width=0.72\textwidth]{lecture_03_figure_04.png}
\caption{The near-horizon causal board: crossed null boundaries, constant-radius hyperbolae in the exterior wedge, and the coordinate relation between \(\rho\) and \(r-2GM\).}
\lectureframe{00:34:35.000}
\end{figure}

\begin{figure}[htbp]
\centering
\begin{tikzpicture}[scale=1.04,>=Latex]
  \draw[->] (-0.3,0) -- (4.6,0) node[right] {\(X\)};
  \draw[->] (0,-3.2) -- (0,3.2) node[above] {\(T\)};
  \draw[very thick,blue!55!black] (0,0) -- (3,3);
  \draw[very thick,blue!55!black] (0,0) -- (3,-3);
  \node[blue!55!black,rotate=45] at (0.72,1.02) {future horizon};
  \draw[thick,domain=-1.45:1.45,samples=80] plot ({0.75*cosh(\x)},{0.75*sinh(\x)});
  \draw[thick,domain=-1.25:1.25,samples=80] plot ({1.25*cosh(\x)},{1.25*sinh(\x)});
  \draw[very thick,domain=-1.05:1.05,samples=80] plot ({2.0*cosh(\x)},{2.0*sinh(\x)});
  \node at (3.64,2.18) {Bob: \(\rho=\mathrm{const.}\)};
  \draw[very thick,red!65!black,->] (1.2,-2.9) -- (1.2,2.9);
  \node[red!65!black] at (1.58,-1.75) {Alice};
  \fill[red!65!black] (1.2,1.2) circle (1.7pt);
  \draw[dashed,->] (1.2,0) -- (2.27,1.07);
  \node[align=left] at (3.18,0.80) {\small pre-crossing\\[-0.2em]\small light signal};
\end{tikzpicture}
\caption{A typeset companion to the board. Bob follows an accelerated constant-\(\rho\) hyperbola; Alice follows a locally inertial line through the future horizon. The null boundary limits Bob's causal access without becoming a local obstacle for Alice.}
\end{figure}

\section{Alice Falls and Bob Hovers}

Let Bob maintain a fixed distance from the horizon while Alice releases his hand and falls. Bob's worldline is a hyperbola and is therefore accelerated. Alice's local worldline is approximately straight. She crosses the null boundary in finite proper time and encounters no local discontinuity.

The causal descriptions are asymmetric. After Alice crosses, light from her cannot reach Bob, while light emitted by Bob can still reach Alice for a time. Bob receives signals emitted by Alice at successively earlier stages of her finite infall. Their arrival times stretch farther and farther apart, their frequency redshifts, and their intensity fades. In Schwarzschild time he never sees the crossing itself; Alice appears to approach and flatten against the horizon asymptotically. Her nose does not flatten while the back of her head remains ordinary: signals from every part are compressed into the same limiting layer.

\begin{classroomqa}
\audiencequestion{Is Alice able to see Bob after crossing because Bob's photons travel inward toward her?\lecturetimestamp{00:42:48}}
\lecturerresponse{Yes. The one-way causal relation is the point: Bob's inward signals can enter, whereas Alice's later outward signals cannot escape back to Bob.}
\end{classroomqa}

\begin{classroomqa}
\audiencequestion{If Alice's feet cross first while her eyes remain outside, does she temporarily lose sight of her own feet?\lecturetimestamp{00:43:03}}
\lecturerresponse{There is no visual blip during an ordinary fall. Light always takes time to travel from feet to eyes; by the time the relevant signal arrives, her eyes have crossed too. Only a violent attempt to hold her head outside while her feet continue inward would create a different and destructive situation.}
\end{classroomqa}

Holding position is the noninertial choice. Bob needs a rocket, just as a person standing on the floor needs the floor's upward force to avoid free fall. The closer Bob hovers to the horizon, the larger the required proper acceleration; in the local Rindler approximation its magnitude is \(a=1/\rho\) in units \(c=1\).\editorialnote{Restoring units gives \(a=c^2/\rho\) for the local Rindler observer. An exactly static Schwarzschild observer has a proper acceleration that diverges as the horizon is approached.}

\begin{classroomqa}
\audiencequestion{Are we accelerating while standing here in the lecture room?\lecturetimestamp{00:45:40}}
\lecturerresponse{Yes, relative to a freely falling frame. The floor supports us; without it, or without a jet pack, we would fall.}
\end{classroomqa}

\begin{classroomqa}
\audiencequestion{Will Bob eventually see Alice within one Planck length of the horizon, and what happens after that?\lecturetimestamp{00:49:24}}
\lecturerresponse{That is where the classical geometric description ceases to be sufficient. The Planck length contains \(\hbar\), so answering what Bob resolves there requires quantum mechanics.}
\end{classroomqa}

\begin{classroomqa}
\audiencequestion{The constant-radius curves seem to meet the horizon at a right-angle-like coordinate corner. Can a derivative be taken there?\lecturetimestamp{00:51:12}}
\lecturerresponse{Something singular happens to this coordinate chart at that point, but not to local physics. A chart that covers the other wedges removes the artificial corner. For the present exterior discussion, the Rindler picture already contains the needed causal facts.}
\end{classroomqa}

It is therefore unhelpful to ask, without naming an observer, whether Alice \emph{really} crosses. Alice's clock reaches the horizon in finite proper time. Bob's received signal sequence never contains the crossing event. Physics asks what each observer measures and how their records can be compared.

\begin{classroomqa}
\audiencequestion{Is the sharp statement that Bob can see only a finite portion of Alice's future, even though his own observation time is unbounded?\lecturetimestamp{00:53:33}}
\lecturerresponse{Yes. Alice assigns a finite proper time to the visible portion of her fall. Bob receives its increasingly delayed and redshifted signals forever in the idealized eternal geometry, but he does not see her continue aging beyond the horizon crossing.}
\end{classroomqa}

\section{Acceleration Horizons Are Already Enough}

Nothing in this causal structure requires a black hole. A rocket that accelerates uniformly forever follows a Rindler hyperbola in flat spacetime and has a null boundary beyond which no signal can catch it. For that observer there is an acceleration horizon, even though an inertial observer regards the same region as ordinary Minkowski space.

\begin{classroomqa}
\audiencequestion{Is this phenomenon specific to black holes, or does it occur for any horizon?\lecturetimestamp{00:55:25}}
\lecturerresponse{The same phenomenon occurs for a uniformly accelerated observer. Such an observer has an acceleration horizon and cannot see events beyond it. Locally, the black-hole horizon is not physically different from that case.}
\end{classroomqa}

This is the equivalence principle at work: understand uniform acceleration first, then recognize the same local structure in a gravitational field. In retrospect, one could have anticipated gravitational horizons from acceleration horizons even before black holes were accepted. Einstein himself did not accept the physical formation of black holes, despite having supplied the theory from which their geometry follows. The horizon is boring for the freely falling observer and consequential for the observer who insists on remaining outside. That is a tension between valid descriptions, not yet a paradox.

\section{Returning to Heat and Entropy}

After the break, put on the thermodynamic hat. In the previous construction, a black hole was filled like a bathtub, one drop at a time, except that each drop was a photon carrying roughly one bit. Each bit increased the horizon area by a quantity of order a Planck area. The exact entropy is
\begin{equation}
S_{\mathrm{BH}}=\frac{k_B A}{4\ell_P^2}.
\end{equation}
The lecture jokes that the factor \(1/4\) could have been avoided had Planck defined his length differently. The serious point is that dimensional analysis fixes the scale but not the numerical coefficient; Hawking's calculation fixes the coefficient.\editorialnote{The Planck length has a standard definition \(\ell_P=\sqrt{G\hbar/c^3}\). The factor \(1/4\) is physical once this convention and the usual thermodynamic entropy are fixed; the remark about redefining it is a classroom joke about conventions.}

Temperature feels elementary because fingers and thermometers register it directly. Entropy often feels more mysterious. The lecture reverses that hierarchy: entropy and energy are treated as the basic variables, while temperature is derived from how much energy must change when entropy changes:
\begin{equation}
dE=T\,dS.
\end{equation}
In the lecture's one-bit units, \(\Delta S=1\), so the minimum energy associated with adding or erasing a bit is called \(T\). Erasing carelessly---with the logical equivalent of a sledgehammer---can dissipate arbitrarily more energy. The definition concerns the least dissipation possible when the missing information is transferred into microscopic environmental degrees of freedom. In conventional units, one binary bit carries \(\Delta S=k_B\ln2\), and reversible erasure dissipates at least \(k_BT\ln2\).\editorialnote{This is the Landauer normalization suppressed by the lecture's choice to measure both entropy and temperature in energy-based natural units.}

\subsection{The one-photon estimate}

Choose a photon whose wavelength is comparable to the Schwarzschild radius:
\begin{equation}
\lambda\sim R_s=\frac{2GM}{c^2}.
\end{equation}
Its energy is, up to the \(2\pi\) convention distinguishing \(h\) from \(\hbar\),
\begin{equation}
\Delta E\sim \frac{\hbar c}{\lambda}.
\end{equation}
That is the energy cost of the added bit, so
\begin{equation}
k_BT_H\sim\frac{\hbar c}{R_s}
\sim\frac{\hbar c^3}{GM}.
\label{eq:l03-hawking-scale}
\end{equation}
The heuristic argument cannot determine factors such as \(2\) and \(\pi\). The exact Hawking result is
\begin{equation}
\boxed{k_BT_H=\frac{\hbar c^3}{8\pi GM}.}
\label{eq:l03-hawking-exact}
\end{equation}

\begin{figure}[htbp]
\centering
\includegraphics[width=0.72\textwidth]{lecture_03_figure_05.png}
\caption{The completed temperature board: \(R_s=2GM/c^2\), the one-bit scale \(\hbar c/R_s\), and the exact \(8\pi\) Hawking normalization.}
\lectureframe{01:07:48.000}
\end{figure}

Two facts are immediate. First, the temperature vanishes with \(\hbar\), so it is a quantum effect. Second,
\begin{equation}
T_H\propto\frac{1}{M}.
\end{equation}
A large black hole is cold. For one solar mass,
\begin{equation}
T_H\approx6\times10^{-8}\ {\rm K},
\end{equation}
which is of the same order as the lecture's \(10^{-8}\,\mathrm{K}\) estimate and far below the cosmic microwave background temperature of about \(2.7\,\mathrm{K}\).

\subsection{Negative heat capacity}

Put such a black hole in a warmer environment. It absorbs energy, its mass increases, and Equation~\eqref{eq:l03-hawking-exact} says that it becomes colder. It then absorbs still more readily. Put it in a colder environment and the reverse runaway begins: it radiates, loses mass, becomes hotter, and radiates faster. Formally,
\begin{equation}
C_{\mathrm{BH}}
=\frac{dE}{dT_H}
=c^2\frac{dM}{dT_H}
<0.
\end{equation}

\begin{classroomqa}
\audiencequestion{Is this what is meant by negative specific heat?\lecturetimestamp{01:14:08}}
\lecturerresponse{Exactly. A self-gravitating system can heat up while losing energy. A star does something analogous: as it radiates, reduced pressure lets it contract, and gravitational compression raises its temperature.}
\end{classroomqa}

An isolated black hole is not thereby unphysical. In today's universe, a stellar hole is cooler than its surroundings and predominantly absorbs. As the universe expands and the background cools, the balance can eventually reverse and evaporation can dominate. The simple equilibrium with an effectively infinite heat bath is unstable.\editorialnote{Boundary conditions matter. A black hole can be stabilized in a finite reflecting cavity or, for appropriate large holes, in anti-de Sitter spacetime; the instability described here is the asymptotically flat, ordinary-bath case.}

\section{The Evaporation Law}

For the scaling derivation, imagine being without a reference book and choose the units in which the remembered equations are shortest:
\begin{equation}
c=\hbar=G=k_B=1.
\end{equation}
These are Planck units. The characteristic scales quoted in the lecture are
\begin{equation}
\ell_P\sim10^{-35}\ {\rm m},
\qquad
t_P\sim10^{-43}\ {\rm s},
\qquad
m_P\sim10^{-8}\ {\rm kg}.
\end{equation}
The Planck mass is roughly a visible dust mote, whereas the Planck length and time are extraordinarily small. Since \(c=1\), \(t_P\) is simply the light-crossing time of \(\ell_P\).

Treat the hole as a black body, meaning a thermal emitter rather than an object that must remain perfectly dark. It can emit electromagnetic radiation and, in principle, other available species including gravitons. The radiated power is proportional to the emitting area because radiation leaves through the surface. The Stefan--Boltzmann scaling is
\begin{equation}
L=-\frac{dE}{dt}\sim A T^4.
\end{equation}
In Planck units \(E=M\), \(R_s\sim M\), \(A\sim R_s^2\sim M^2\), and \(T\sim R_s^{-1}\sim M^{-1}\). Therefore
\begin{align}
-\frac{dM}{dt}
&\sim R_s^2\left(\frac{1}{R_s}\right)^4
\sim\frac{1}{R_s^2}
\sim\frac{1}{M^2},\\
-M^2\frac{dM}{dt}&\sim1,\\
-\frac{d(M^3)}{dt}&\sim1.
\label{eq:l03-evap-scaling}
\end{align}
Ignoring order-one coefficients, the lifetime is
\begin{equation}
\boxed{t_{\mathrm{evap}}\sim M^3.}
\label{eq:l03-lifetime}
\end{equation}

\begin{figure}[htbp]
\centering
\includegraphics[width=0.72\textwidth]{lecture_03_figure_06.png}
\caption{The unobstructed evaporation board: luminosity gives \(-dM/dt\sim1/M^2\), hence a uniform decrease of \(M^3\) and a lifetime proportional to the initial mass cubed.}
\lectureframe{01:23:20.000}
\end{figure}

The schematic Stefan--Boltzmann treatment suppresses greybody factors and the number of particle species. For a neutral, nonrotating hole in the simplest massless-emission approximation, restoring constants gives
\begin{equation}
t_{\mathrm{evap}}
=\frac{5120\pi G^2M_0^3}{\hbar c^4}.
\label{eq:l03-exact-lifetime}
\end{equation}
This coefficient is an editorial completion; the lecture is deriving only the robust \(M_0^3\) dependence.

\subsection{Solar, mountain, and kilogram scales}

The Sun has mass \(M_\odot\sim10^{30}\,\mathrm{kg}\sim10^{38}m_P\). Cubing and multiplying by one Planck time gives the lecture's deliberately rough chain
\begin{equation}
t_{\mathrm{evap}}\sim10^{114}t_P
\sim10^{71}\ {\rm s}
\sim10^{64}\ {\rm yr}.
\end{equation}
Either estimate is fantastically longer than the roughly \(10^{10}\)-year age of the universe. Restoring the coefficient in Equation~\eqref{eq:l03-exact-lifetime} raises the standard solar-mass estimate to about \(2\times10^{67}\) years.\editorialnote{The difference is expected: the blackboard calculation drops all numerical constants and rounds the Planck scales by powers of ten. Its purpose is the mass-cubed scaling, not precision lifetime astronomy.}

The class next tries a \emph{mountain-mass} hole. It settles on the rough classroom value \(M\sim10^{10}\,\mathrm{kg}\), hence \(M\sim10^{18}m_P\), and obtains
\begin{equation}
t_{\mathrm{evap}}\sim10^{54}t_P
\sim10^{11}\ {\rm s}
\sim10^4\ {\rm yr}.
\end{equation}
This preserves the arithmetic actually carried through at the board. The verbal geometry used to motivate that mass---roughly \(100\,\mathrm{km}^3\) of rock---would in fact be much heavier, and the exact lifetime coefficient would also change the result substantially.\editorialnote{At ordinary rock density, \(100\,\mathrm{km}^3\) is of order \(10^{14}\,\mathrm{kg}\), not \(10^{10}\,\mathrm{kg}\). For \(10^{10}\,\mathrm{kg}\), Equation~\eqref{eq:l03-exact-lifetime} gives a lifetime of order \(10^6\) years. These corrections do not alter the lecture's lesson that lifetime falls as the cube of mass.}

\begin{classroomqa}
\audiencequestion{Would a black hole with a mass of only a kilogram be extremely small?\lecturetimestamp{01:29:14}}
\lecturerresponse{Yes. Since radius is proportional to mass and a solar mass corresponds to a few kilometres, a one-kilogram Schwarzschild radius is about \(10^{-27}\,\mathrm{m}\). Such a hole would also be extremely hot and short-lived.}
\end{classroomqa}

\begin{classroomqa}
\audiencequestion{Do these evaporation estimates assume zero ambient temperature?\lecturetimestamp{01:29:53}}
\lecturerresponse{They apply when the surroundings are cooler than the hole, so that net energy flows outward. A hotter environment instead feeds the black hole. The mass whose Hawking temperature matches the present background is left as an exercise.}
\end{classroomqa}

For orientation, setting \(T_H=2.7\,\mathrm{K}\) gives \(M\) of order \(10^{22}\,\mathrm{kg}\) and a Schwarzschild radius of order \(10^{-4}\,\mathrm{m}\).\editorialnote{This evaluates the exercise posed in the lecture; it is not part of the spoken derivation.}

\section{What Must Be Under the Metric?}

\begin{classroomqa}
\audiencequestion{What mechanism makes the photons come out of the black hole?\lecturetimestamp{01:30:53}}
\lecturerresponse{Before giving a microscopic mechanism for the radiation, one must identify the degrees of freedom responsible for the entropy. The thermodynamic argument constrains what must exist without yet displaying those variables explicitly.}
\end{classroomqa}

The analogy is fluid mechanics. Hydrodynamics describes smooth density and velocity fields without following every molecule. Entropy nevertheless reveals that many hidden microscopic configurations underlie the same smooth state. In the same way, classical general relativity describes a smooth metric but does not expose the numerous microscopic degrees of freedom counted by black-hole entropy. The area law and Hawking radiation are clues that Einstein's theory is an effective macroscopic description.

\begin{classroomqa}
\audiencequestion{If the black hole radiates energy, is that the same as saying that it radiates information, so that information swallowed during growth returns during evaporation?\lecturetimestamp{01:33:27}}
\lecturerresponse{That is the requirement: if the complete evolution conserves information, the outgoing radiation must eventually distinguish the possible states that formed the hole. This must still be reconciled with the causal picture in which matter crosses the horizon.}
\end{classroomqa}

Energy flux by itself is not enough to specify information. A perfectly thermal density matrix can carry energy while forgetting the detailed initial state. In a unitary description, the required information is encoded in fine correlations among the emitted quanta.\editorialnote{This distinction makes explicit the modern form of the information problem. It sharpens rather than changes the question raised in the lecture.}

\section{Information Conservation as a Structural Principle}

\begin{classroomqa}
\audiencequestion{Is conservation of information truly a fundamental principle?\lecturetimestamp{01:34:13}}
\lecturerresponse{It is among the most fundamental. If distinct initial states could evolve into exactly the same final state, the laws would erase their own past. The usual conservation laws rely not only on symmetries but also on a reversible Hamiltonian or unitary form of evolution.}
\end{classroomqa}

Imagine, as the lecture does, a fictional fundamental law that simply makes every sliding eraser come to rest. Different initial velocities all end in the same state. The final eraser no longer distinguishes its histories: information has been destroyed. Momentum and mechanical energy also disappear. Real friction does not do this at the microscopic level; momentum and energy pass into the table, air, and thermal degrees of freedom.

Even a perfectly uniform infinite tabletop would possess translation symmetry, yet the fictitious damping law could still make the eraser stop. Translation symmetry alone is therefore not enough to derive momentum conservation from such a non-Hamiltonian rule. Noether's theorem assumes an action principle and the associated reversible dynamics. In the language used here, symmetry plus information-preserving dynamics gives the conservation law.

\begin{classroomqa}
\audiencequestion{Is information conservation more fundamental than causality?\lecturetimestamp{01:37:00}}
\lecturerresponse{They are different principles at a comparable foundational level. Causality must first be stated precisely---for example, that signals cannot propagate backward outside the permitted causal structure---but neither principle can simply replace the other.}
\end{classroomqa}

\begin{classroomqa}
\audiencequestion{Is the phase-space area preservation discussed in Hamiltonian mechanics the classical version of this idea?\lecturetimestamp{01:37:38}}
\lecturerresponse{Yes. Liouville's theorem says Hamiltonian flow preserves phase-space volume. A cloud of possible initial states cannot collapse into a smaller volume and erase distinctions. In quantum mechanics the corresponding statement is unitarity: inner products, and hence distinctions between orthogonal states, are preserved.}
\end{classroomqa}

Classically,
\begin{equation}
\frac{d}{dt}\,\mathrm{Vol}_{\Gamma}=0,
\end{equation}
for a phase-space region transported by Hamiltonian flow. Quantum mechanically,
\begin{equation}
\begin{aligned}
\ket{\psi(t)}&=U(t)\ket{\psi(0)},\\
U^\dagger U&=I,\\
\braket{\psi(t)}{\phi(t)}
&=\braket{\psi(0)}{\phi(0)}.
\end{aligned}
\end{equation}
These are precise forms of the claim that microscopic distinctions survive time evolution.

\section{Which Way Does Entropy Depend?}

The final exchange returns to the information-theoretic entropy. Given probabilities \(p_i\) for microscopic alternatives,
\begin{equation}
S=-k_B\sum_i p_i\ln p_i
\label{eq:l03-shannon}
\end{equation}
in thermodynamic units, or without \(k_B\) for dimensionless Shannon entropy.

\begin{figure}[htbp]
\centering
\includegraphics[width=0.72\textwidth]{lecture_03_figure_07.png}
\caption{The closing board formula \(S=-\sum_i p_i\log p_i\), written while distinguishing a probability distribution from the single entropy value computed from it.}
\lectureframe{01:40:32.000}
\end{figure}

\begin{classroomqa}
\audiencequestion{Does maximizing the entropy define the probability distribution, rather than the probabilities defining the entropy?\lecturetimestamp{01:38:33}}
\lecturerresponse{A probability distribution determines one entropy value through Equation~\eqref{eq:l03-shannon}; one number cannot reconstruct all the \(p_i\). Maximizing entropy can select a distribution only after normalization and other constraints, such as fixed mean energy, have also been supplied.}
\end{classroomqa}

That distinction ends the lecture at the right frontier. Smooth geometry has led to a horizon; the horizon has led to thermodynamics; thermodynamics has exposed hidden degrees of freedom; and information conservation now demands that those degrees of freedom evolve without erasing the distinctions that entered the hole. The mechanism has not yet been supplied. The geometry and the bookkeeping have instead made clear what any successful microscopic theory must explain.
